% Generated draft: requires mathematical and source review.
\documentclass[11pt]{article}
\usepackage[margin=25mm]{geometry}
\usepackage{fontspec}
\setmainfont{Latin Modern Roman}
\usepackage{amsmath,amssymb,mathtools,unicode-math}
\setmathfont{Latin Modern Math}
\usepackage{enumitem,needspace,xurl,hyperref}
\setlength{\parindent}{0pt}
\setlength{\parskip}{5pt}
\setlength{\emergencystretch}{3em}
\newcommand{\N}{\mathbb{N}}
\newcommand{\Z}{\mathbb{Z}}
\newcommand{\Q}{\mathbb{Q}}
\newcommand{\R}{\mathbb{R}}
\newcommand{\C}{\mathbb{C}}
\newcommand{\F}{\mathbb{F}}
\newcommand{\A}{\mathbb{A}}
\newcommand{\PP}{\mathbb{P}}
\newcommand{\E}{\mathbb{E}}
\newcommand{\eps}{\varepsilon}
\newcommand{\dd}{\,\mathrm{d}}
\begin{document}
\section*{UnsolvedMath ID: 30001477}
\textbf{Source problem code:} catalogue entry
\subsection{Definitions and mathematical statement}
The normalized source asks for a classification of mod-two area-minimizing surfaces in the Tri, Tetra, and Hexa Cosms, regarded as closed orientable flat three-manifolds. The record does not specify their flat metrics, lattice parameters, or the precise admissible class of surfaces. The following gives a homological formulation conditional on the intended meaning of minimization.

Let \((M,g)\) be one of these manifolds equipped with a fixed flat metric. Write \(\mathcal Z_2(M;\Z/2\Z)\) for the integral two-dimensional cycles with coefficients modulo two, and write \(\mathbf M_g(T)\) for the mass, or area, of such a cycle \(T\). For \(\alpha\in H_2(M;\Z/2\Z)\), define
\[
\mathcal A_g(\alpha)=\left\{T\in\mathcal Z_2(M;\Z/2\Z):[T]=\alpha,\quad \mathbf M_g(T)=\inf_{\substack{S\in\mathcal Z_2(M;\Z/2\Z)\\[S]=\alpha}}\mathbf M_g(S)\right\}.
\]
Under this interpretation, the classification problem is to determine \(\mathcal A_g(\alpha)\) for each homology class \(\alpha\), for each of the three named manifolds. Whether the intended problem concerns a particular metric or all flat metrics remains unspecified. The use of closed cycles, minimization within a homology class, and the inclusion of every homology class require confirmation from the source context.

To make the subsidiary planarity question precise, call a smooth surface planar if it is totally geodesic, equivalently if its lifts to Euclidean covering coordinates are locally contained in affine planes. The phrase ``planar-minimizer conclusion'' could ask whether every minimizing representative is planar, or whether each relevant class admits a planar minimizing representative. These are different quantified assertions, and the supplied record does not resolve which is intended. It also does not specify whether disconnected surfaces or generalized cycles are part of the intended classification.

The original formulation concerns compact flat three-manifolds other than the Di and Didi Cosms, whereas the normalized formulation restricts the target to the orientable Tri, Tetra, and Hexa Cosms and treats the three-torus, Di Cosm, and Didi Cosm as covered by known planar-minimizer results. The broader wording could also encompass nonorientable manifolds; that scope difference is unresolved here.
\subsection{Short English statement}
Classify mod-two area minimizers in the Tri, Tetra, and Hexa Cosms and clarify the corresponding planarity question. The intended metrics, comparison class, quantifiers for planarity, and relationship to the broader original formulation require source review.
\subsection{Sources}
\begin{itemize}
\item Supplied local UnsolvedMath registry record, permanent ID \(30001477\), source problem code \texttt{catalogue entry}, title: ``Mod-Two Area Minimizers in Compact Flat Three-Manifolds.'' The record assigns the problem to 2010 and supplies no named proposer.
\item \emph{Progress in Surface Theory} (2010), Oberwolfach Reports, pp.~1225--1318, DOI \texttt{10.4171/owr/2010/21}. Supplied report reference and literature-source link: \url{https://doi.org/10.4171/owr/2010/21}. The supplied background describes Oberwolfach Reports as documenting mathematical workshops and their open-problem discussions; no problem-specific page or contributor is identified.
\item Supplied dated literature curation, checked \texttt{2026-08-21}: the status is recorded as ``open,'' and the assessment states that the report proves planar cases for two compact flat three-manifolds and explicitly leaves the remaining cases open. The normalized statement separately describes planar-minimizer results as known for the three-torus, Di Cosm, and Didi Cosm. These are supplied source and literature claims, not independently verified mathematical or current-status conclusions.
\item The registry labels the normalized statement \texttt{corrected\_verified} and reports that its mathematical content was checked against the official Oberwolfach report. This verification claim is supplied metadata; no independent source verification has been performed for this draft.
\end{itemize}
\end{document}
